\item Un electrón está en la $n$-ésima órbita de Bohr del átomo de hidrógeno.
\begin{enumerate}
    \item Demuestre que el periodo orbital del electrón es $T = n^3 t_0$ y determine el valor numérico de $t_0$.
    \item En promedio, un electrón permanece en la órbita $n = 2$ durante aproximadamente $10\text{ }\mu\text{s}$ antes de saltar hacia la órbita $n = 1$ (estado fundamental). ¿Cuántas revoluciones realiza el electrón en este estado excitado?
    \item Defina el periodo de una revolución como un ``año electrón'', análogo a un año Tierra. Explique si debe considerar al electrón en la órbita $n = 2$ como ``viviendo por un largo tiempo''.
\end{enumerate}